% !TeX encoding = UTF-8
% !TeX spellcheck = de_DE
% Datei: 00_vorwort.tex
% Ein Vorwort, das die analytische Schärfe der Theorie demonstriert

\chapter*{Vorwort: Eine Theorie als seismographisches Instrument}
\addcontentsline{toc}{chapter}{Vorwort}
\label{chap:vorwort}

\thispagestyle{empty}

Wozu dient eine Gesellschaftstheorie? Sie kann ein utopischer Entwurf, ein historisches Erklärungsmodell oder eine normative Richtschnur sein. Die
vorliegende Theorie versteht sich vor allem als etwas anderes: als ein \textbf{seismographisches Instrument}. Ihr Wert liegt weniger in den
konkreten Institutionen, die sie entwirft, als in der erbarmungslosen Klarheit, mit der sie bestehende politische Ordnungen durchleuchtet und deren
verborgene Widersprüche aufdeckt.

Ein einfacher Test genügt, um diese analytische Schärfe zu demonstrieren. Nehmen wir eine scheinbar selbstverständliche, friedlich etablierte
Struktur unserer politischen Wirklichkeit: das mehrstufige \textbf{föderale System} aus Kommunen, Bundesländern, Nationalstaat und supranationaler
Union (wie der EU). Es erscheint uns als ausgewogener Kompromiss zwischen lokaler Autonomie und übergeordneter Handlungsfähigkeit, zwischen
Demokratie und Effizienz.

Wenden wir nun die einfache Frage an, die das Herzstück dieser Theorie bildet: \textbf{„Ist diese Struktur \emph{wirklich} demokratisch?“}

Die Antwort, die sich aus den folgenden Axiomen zwingend ergibt, ist verstörend eindeutig: \textbf{Nein.} Sie ist es strukturell nicht und kann es
nicht sein. Denn bereits die zweite Stufe dieses Systems – die repräsentative Demokratie auf Landes- oder Bundesebene – bricht mit einem
fundamentalen Prinzip: Sie konzentriert notwendigerweise politische Macht in dauerhaften Organisationen (Parteien, Parlamentsfraktionen,
Regierungsapparaten) und entzieht sie der direkten Kontrolle der Einzelnen. Die Wahl allein vier Jahre genügt nicht, um aus einem
Repräsentationssystem eine „wirkliche Demokratie“ zu machen; sie institutionalisiert vielmehr eine periodisch erneuerte
\textbf{Delegation von Souveränität}.

Gehen wir weiter: Was ist eine „Landesregierung“ oder eine „Bundesregierung“ anderes als eine – durch Wahlen lediglich legitimierte –
\textbf{Machtkonzentration}, die für Jahre über das Leben von Millionen entscheidet? Was ist die Europäische Kommission anderes als eine exekutive
Elite, deren demokratische Rückkopplung über ein mehrstufiges, stark gefiltertes System von Delegationen so verdünnt ist, dass sie kaum noch spürbar
ist?

Die hier entwickelte Theorie nennt diese Dinge beim Namen. Sie postuliert in ihrer Radikalität, dass jede stabile Organisation, deren Zweck die
Ausübung politischer Macht ist, per Definition undemokratisch ist (Axiom 2). Folgerichtig kann ein föderaler Staat mit seinen Regierungsebenen aus
dieser Perspektive nur als eine \textbf{Hierarchie von Machtkonzentrationen} verstanden werden, die sich über die eigentlichen Träger der
Souveränität – die Individuen in ihren Lebenswelten – geschichtet hat.

Die scheinbare Stabilität und Effizienz dieses Systems erkaufen wir mit einem systematischen \textbf{Demokratiedefizit}. Dieses Defizit ist kein
bedauerlicher Unfall, sondern das konstitutive Merkmal der repräsentativ-föderalen Ordnung.

Dieses Buch führt diese radikale Kritik nicht als destruktive Geste aus. Im Gegenteil: Sie ist der Ausgangspunkt für den konstruktiven Versuch, zu
ergründen, wie eine Gesellschaft beschaffen sein müsste, die den Namen „Demokratie“ ohne Vorbehalt verdient. Die fünf Axiome, die im ersten Kapitel
entfaltet werden, sind der Versuch, die Grundlagen für eine solche Gesellschaft zu legen. Sie sind kompromisslos, vielleicht sogar utopisch. Doch
ihr größter Wert liegt vielleicht gerade in dieser Kompromisslosigkeit: Sie zwingen uns, jede einzelne Institution, von der wir umgeben sind, erneut
zu befragen.

Kann es Demokratie in Millionenverbänden geben? Muss eine demokratische Gesellschaft notwendig klein oder dezentral sein? Kann wissenschaftliche
Rationalität mit Volkssouveränität versöhnt werden, ohne in Technokratie umzuschlagen? Diese Theorie gibt provokante Antworten, die nicht gefallen
werden. Aber sie zwingt zur Stellungnahme.

Dieses Buch ist somit eine Einladung. Eine Einladung, die selbstverständlich gewordenen Grundlagen unseres Zusammenlebens nicht als naturgegeben
hinzunehmen, sondern sie durch das Prisma einer konsequenten Idee zu betrachten: der Idee, dass die Herrschaft des Volkes wörtlich zu nehmen ist.
Das Ergebnis wird unbequem sein. Es wird zeigen, wie weit wir von einer „wirklichen Demokratie“ entfernt leben. Doch nur wer das Ausmaß der
Entfernung kennt, kann überhaupt einen Weg einschlagen.

\vspace{1cm}
\hfill Michael Czybor
\hfill \textit{\today}